\documentclass{article}
\usepackage[utf8]{inputenc}
\usepackage{a4wide}

\begin{document}

\section*{เลขที่หายไป}

กำหนดอาร์เรย์ของเลขจำนวนเต็ม\texttt{nums}ขนาด\texttt{N}ที่เรียงลำดับเพิ่มขึ้นต่อเนื่องและไม่มีเลขซ้ำ

และกำหนดจำนวนเต็ม\texttt{K}ให้

ให้หาเลขที่หายไปตัวที่\texttt{Kth}ในชุดเลข\texttt{nums}นี้ (เริ่มจากซ้ายมือ)\\



\textbf{ตัวอย่าง}



{\def\LTcaptype{none} % do not increment counter

\begin{longtable}[]{@{}

  >{\raggedright\arraybackslash}p{(\linewidth - 4\tabcolsep) * \real{0.3333}}

  >{\raggedright\arraybackslash}p{(\linewidth - 4\tabcolsep) * \real{0.3333}}

  >{\raggedright\arraybackslash}p{(\linewidth - 4\tabcolsep) * \real{0.3333}}@{}}

\toprule\noalign{}

\begin{minipage}[b]{\linewidth}\raggedright

ข้อมูลเข้า

\end{minipage} & \begin{minipage}[b]{\linewidth}\raggedright

ข้อมูลออก

\end{minipage} & \begin{minipage}[b]{\linewidth}\raggedright

คำอธิบาย

\end{minipage} \\

\midrule\noalign{}

\endhead

\bottomrule\noalign{}

\endlastfoot

\begin{minipage}[t]{\linewidth}\raggedright

4 3\\

4 7 9 10\strut

\end{minipage} & 8 & 5,6,\textbf{8} \\

\end{longtable}

}



\subsection*{Input}
บรรทัดที่ 1 จำนวนเต็ม\texttt{N}เป็นขนาดของอาร์เรย์

และตามด้วย~เลขจำนวนเต็ม\texttt{K}เว้นด้วยช่องว่าง



บรรทัดที่ 2{จำนวนเต็ม}\texttt{nums}{ทั้งหมด}\texttt{N}{ตัวเว้นด้วยช่องว่าง}



\subsection*{Output}
{เลขที่หายไปตัวที่}\texttt{Kth}{ในชุดเลข}\texttt{nums}{นี้}\\



\end{document}
